% Shared IEEE conference formatting. Keep the class, margins and body size intact.
\documentclass[conference,letterpaper]{IEEEtran}
\usepackage[T1]{fontenc}
\usepackage{amsmath,amssymb}
\usepackage{cite}
\usepackage{booktabs,array}
\usepackage{balance}
\usepackage[hidelinks]{hyperref}
\urlstyle{same}
\hypersetup{pdfdisplaydoctitle=true}

\hypersetup{
  pdftitle={Fostering Cellular Agriculture},
  pdfauthor={Santiago Restrepo Ruiz},
  pdfsubject={One-page synthesis of Working Paper v2.1: Financing Cellular-Agriculture Scale-Up under Correlated Technical Risk},
  pdfkeywords={cellular agriculture, financial engineering, milestone participation,
    portfolio risk, loss allocation, pricing, funding liquidity}
}

\title{Fostering Cellular Agriculture}
\author{\IEEEauthorblockN{Santiago Restrepo Ruiz}
\IEEEauthorblockA{Published by Waajacu Open Source Foundation\\
\href{mailto:Santiago.Restrepo.Ruiz@gmail.com}{Santiago.Restrepo.Ruiz@gmail.com}\quad
\url{https://waajacu.com/}\\[3pt]
\textit{One-page synthesis of Working Paper v2.1, 1 September 2026}}}
\date{}

\begin{document}
\maketitle

\begin{abstract}
Fostering Cellular Agriculture proposes an open financial standard for
research-to-industry funding. Its Multi-Project Milestone Participation connects
staged capital to enforceable project cash rights and recoveries. A common
claim interface, explicit portfolio dependence, contractual waterfalls and
funding controls support term design against investor return, loss and
liquidity constraints.
\end{abstract}

\begin{IEEEkeywords}
Cellular agriculture, financial engineering, milestone finance.
\end{IEEEkeywords}

\balance
\section{Claim Design and Milestone Funding}
Cellular-agriculture scale-up links uncertain process economics to staged
investment~\cite{humbird2021scale}. Research-backed obligations provide a
precedent for pooling development assets~\cite{fernandez2012rbo}; this
structure begins with each project's assignable cash rights.

Each project specifies dated capital needs, verified milestones, success cash,
recovery priority, principal and shared risk factors. Commercial, licensing,
royalty, capacity, offtake or exit receipts trace to contractual rights and
non-overlapping cash-source budgets.

Investors hold a limited-recourse participation in a fixed pool. Milestones
release callable funding; failure can stop later draws. Receipts and recoveries
pass through pro rata after pool cost, without a coupon independent of project
cash. Acquisition cost, asset principal and issued principal remain distinct;
unresolved principal remains exposure, not cash or realized loss.

\section{Portfolio Risk and Loss Allocation}
Joint cash-flow scenarios encode shared biology, scale-up, supplier, regulatory
and buyer risks. Dependence is explicit~\cite{hull2019longshots}; bounded
physical state probabilities expose uncertainty. Expected shortfall measures
loss and net present value (NPV) downside tails. Diversification compares pooled and standalone losses
under the same probability measure.

A funded variant divides claims into first-loss and priority layers.
Principal cash and unused reserve repay senior claims first; non-principal
cash follows a capped priority waterfall, then the residual claim. Layers
allocate the issued-principal cash shortfall, separately from asset writeoffs
and continuing exposure. Returned reserve is capital, not profit.

Alternatively, a capped partial-credit guarantee covers a share of the
untranched core's resolved asset loss, excluding continuing exposure.
Collateral and joint project/provider credit stress determine delivered cash.
Tranching redistributes pool cash; guarantees introduce external transfers.

\section{Pricing and Term Selection}
For net investor cash \(CF_s(t)\) in state \(s\), month \(t\), supplied annual
hurdle \(h\) and probability set \(\mathcal P\), robust return requires
\begin{equation*}
\min_{p\in\mathcal P}\sum_s p_s\sum_t
\frac{CF_s(t)}{(1+h)^{t/12}}\;\geq\;0.
\end{equation*}
Success-participation and first-loss terms are screened jointly against return,
principal-shortfall tails, impairment probability, weighted-average life and
junior-concession limits. Each adverse endpoint retains its probability witness.

For a fixed claim, compare the investor's price ceiling with the issuer's
funding floor, including costs and support. Discounts cannot repair fixed
principal-risk breaches. Compare guarantee premium capacity with provider
claims, expenses, liquidity/capital costs and required return; any gap requires
external support. These are conditional value tests, not market quotes.

\section{Funding Integrity and Execution}
Full prefunding trades capital-call risk for idle-cash carry. A proposed
callable-capital or warehouse alternative specifies borrowing-base eligibility,
collateral, maturity, repayment priority and provider capacity. Decisions use
only observed information. Settled funding covers uses before later project
receipts; undrawn commitments are not cash. Funding gaps invalidate subsequent
assumed project cash paths.

Implementation reconciles cash, principal and waterfalls in every state.
Underwriting requires authenticated contracts, milestone and recovery histories,
joint-risk evidence, observed investor prices or hurdles, and executable funding
terms. The output identifies what investors own, how they are paid, who bears
each risk, and whether the complete structure satisfies its constraints.

\bibliographystyle{IEEEtran}
\bibliography{references}
\end{document}
